\section*{ID8499: DIMENSIONAL REDUCTION CORRELATION Canonical Statutory Formula}
\paragraph{Metadata.}
PiID: 8340800535598 | RTEID: RTE:P33003.5044:T4.1141 | UUID: ec1b6cfd-07f8-5a32-91ca-3153c67503a5

\paragraph{Canonical Statement.}
[ID 1173] DIMENSIONAL REDUCTION CORRELATION Canonical Statutory Formula: \$\textbackslash{}langle O\_1(0) O\_2(r) \textbackslash{}rangle \textbackslash{}propto r\textasciicircum{}\{-(N-1)\}\$ Formal Technical Dissertation: A verification metric for criticality. It predicts that at the Trawin Zero Point, an \$N\$-dimensional system effectively computes on an \$(N-1)\$-dimensional boundary (the null surface), creating a specific decay signature in correlation functions. Source: [21]

\paragraph{Canonical Equation.}
\[
\langle O_1(0) O_2(r) \rangle \propto r^{-(N-1)}
\]

\paragraph{Dictionary.}
DIMENSIONAL REDUCTION CORRELATION Canonical Statutory Formula: ⟨ O\_1(0) O\_2(r) ⟩ ∝ r\textasciicircum{}-(N-1) Formal Technical Dissertation: A verification metric for criticality.

\paragraph{Encyclopedia.}
DIMENSIONAL REDUCTION CORRELATION Canonical Statutory Formula is an expectation / trace over state, written ⟨ O\_1(0) O\_2(r) ⟩ ∝ r\textasciicircum{}-(N-1). It is computed as an expectation value or trace over a quantum state, producing an observable. It is built from O\_1, O\_2, r, N. Within the Trawin Zero Point registry it serves as one element of the framework's mathematical narrative.
